\chapter{Coding Course}

\section{Goal Of This Course}

This chapter teaches the minimum coding knowledge needed to read, run, and
debug this package. It does not try to teach all of C++ or all of ROS 2. It
teaches the ideas that appear in this repository:

\begin{itemize}
  \item source files and header files,
  \item variables, functions, structs, and classes,
  \item loops and conditions,
  \item ROS 2 nodes, topics, messages, publishers, and subscribers,
  \item CMake and package metadata,
  \item practical debugging habits.
\end{itemize}

\section{What A Program Is}

A program is a set of instructions for the computer. Source code is the human
readable text that describes those instructions. A compiler translates C++
source code into machine code that the computer can run.

This package is a ROS 2 package named \repo. ROS 2 is a robotics framework. It
lets separate programs communicate by sending messages on named channels called
topics.

\section{Files, Folders, And Builds}

The important hand-written code is under \file{src/}. Generated files under
\file{build/} and \file{install/} are build output. New members should usually
read and edit \file{src/}, not generated output.

\begin{longtable}{p{0.28\textwidth}p{0.62\textwidth}}
\toprule
Folder or file & Beginner explanation \\
\midrule
\file{src/src/*.cpp} & C++ implementation files. These contain function
bodies: what the code actually does. \\
\file{src/include/auto_aim/*.hpp} & C++ header files. These declare the
types and functions other files can use. \\
\file{src/msg/*.msg} & ROS message definitions. ROS generates C++ code from
these during the build. \\
\file{src/config/*.yaml} & Parameter values loaded at runtime. This is where
robot-specific tuning belongs. \\
\file{src/launch/*.py} & Launch scripts that start the node with parameters. \\
\file{src/CMakeLists.txt} & Build instructions for CMake and colcon. \\
\file{src/package.xml} & ROS package metadata and dependencies. \\
\bottomrule
\end{longtable}

Typical build and run commands from the workspace root:

\begin{lstlisting}
colcon build --packages-select auto_aim
source install/setup.bash
ros2 launch auto_aim auto_aim_realsense_16.launch.py
\end{lstlisting}

If you edit C++ or message files, rebuild. If you only edit YAML parameters,
you usually only need to relaunch.

\section{C++ Source And Header Files}

C++ projects often split code into headers and source files.

\begin{itemize}
  \item A header file, such as \file{tracker.hpp}, says what a class or
  function looks like.
  \item A source file, such as \file{tracker.cpp}, contains the implementation.
\end{itemize}

Example from the tracker header:

\begin{lstlisting}[language=C++]
class Tracker
{
public:
  explicit Tracker(const TrackerConfig & cfg);
  void update(const std::vector<ArmorDetection> & detections, double dt);
  AimResult computeAim(double current_yaw, double current_pitch,
                       double override_pred_lead_s = -1.0) const;
};
\end{lstlisting}

This says a \code{Tracker} has a constructor, an \code{update} function, and a
\code{computeAim} function. The detailed math lives in \file{tracker.cpp}.

\section{Includes}

At the top of a C++ file, \texttt{\#include} brings in declarations from another
file or library:

\begin{lstlisting}[language=C++]
#include <rclcpp/rclcpp.hpp>
#include "auto_aim/tracker.hpp"
\end{lstlisting}

Angle brackets usually mean an external library or system header. Quotes
usually mean a header from this package.

\section{Namespaces}

The code is wrapped in:

\begin{lstlisting}[language=C++]
namespace auto_aim
{
  // code here
}
\end{lstlisting}

A namespace is a named container that prevents name collisions. The class is
\code{auto_aim::Tracker}, not just \code{Tracker}, when referred to from
outside the namespace.

\section{Variables And Types}

A variable stores a value. A type tells C++ what kind of value it stores.

\begin{lstlisting}[language=C++]
double bullet_speed = 25.0;
bool tracking = false;
int confirm_frames = 3;
std::string class_id = "0";
\end{lstlisting}

Common types in this repo:

\begin{longtable}{p{0.24\textwidth}p{0.66\textwidth}}
\toprule
Type & Meaning \\
\midrule
\code{double} & Floating point number, used for math. \\
\code{float} & Smaller floating point number, often used in ROS messages. \\
\code{int} & Whole number. \\
\code{bool} & True or false. \\
\code{std::string} & Text. \\
\code{std::vector<T>} & Resizable list of values of type \code{T}. \\
\code{std::array<T,N>} & Fixed-size list with exactly \code{N} values. \\
\code{Eigen::VectorXd} & Math vector from the Eigen library. \\
\code{cv::Mat} & Matrix/image container from OpenCV. \\
\bottomrule
\end{longtable}

\section{Structs}

A struct groups related values under one name. For example:

\begin{lstlisting}[language=C++]
struct ArmorDetection
{
  double x, y, z;
  double yaw;
  std::string class_id;
  double confidence;
};
\end{lstlisting}

Instead of passing six separate values everywhere, the code passes one
\code{ArmorDetection}.

Important structs in this package:

\begin{itemize}
  \item \code{TrackerConfig}: all tracker, ballistic, and gate parameters.
  \item \code{ArmorDetection}: one transformed detection in odom frame.
  \item \code{AimResult}: output from the aim planner.
  \item \code{DebugFrame}: internal per-frame debug data before publishing.
  \item \code{PnPResult}: output from the PnP solver.
\end{itemize}

\section{Functions}

A function is reusable code with a name. It may take inputs and return an
output.

\begin{lstlisting}[language=C++]
double Tracker::targetRange() const
{
  return std::sqrt(x_(0)*x_(0) + x_(2)*x_(2) + x_(4)*x_(4));
}
\end{lstlisting}

This function returns a \code{double}. The \code{const} at the end means it
does not modify the \code{Tracker} object.

\section{Classes}

A class groups data and functions together. In this repo, each major pipeline
piece is a class or stateless helper:

\begin{longtable}{p{0.24\textwidth}p{0.66\textwidth}}
\toprule
Class or helper & Responsibility \\
\midrule
\code{AutoAimNode} & ROS wiring: parameters, subscriptions, callbacks,
publishers. \\
\code{PnPSolver} & 2D box to 3D armor pose. \\
\code{FrameTransformer} & Camera frame to odom frame. \\
\code{Tracker} & EKF, target association, aiming, and ballistics. \\
\code{FireGate} & Fire/hold decision and reason. \\
\code{DetectionAdapter} & Convert detector messages into candidates. \\
\code{DebugPublisher} & Publish structured debug messages. \\
\code{ConfigValidator} & Validate parameters at startup. \\
\bottomrule
\end{longtable}

\section{Public And Private}

Classes often have \code{public} and \code{private} sections.

\begin{itemize}
  \item \code{public}: other code is allowed to call these functions or read
  these members.
  \item \code{private}: internal implementation details. Other code should not
  depend on them.
\end{itemize}

For example, other code can call \code{tracker_->update(...)}, but it cannot
directly change \code{x_}, the EKF state. That keeps the tracker responsible
for its own consistency.

\section{Pointers And Unique Ownership}

Some objects are stored through pointers:

\begin{lstlisting}[language=C++]
std::unique_ptr<Tracker> tracker_;
tracker_ = std::make_unique<Tracker>(cfg);
\end{lstlisting}

A pointer stores the location of an object. \code{std::unique_ptr} means one
owner is responsible for that object. When the owner is destroyed, the object
is automatically cleaned up.

To call a function through a pointer, use \code{->}:

\begin{lstlisting}[language=C++]
tracker_->update(armors, dt);
\end{lstlisting}

\section{References And \texorpdfstring{\code{const}}{const}}

You will often see parameters like:

\begin{lstlisting}[language=C++]
void update(const std::vector<ArmorDetection> & detections, double dt);
\end{lstlisting}

\code{&} means pass by reference. The function can use the original object
without copying the whole vector. \code{const} means the function promises not
to modify that vector.

This is common for performance and safety.

\section{Conditions}

\code{if} chooses between code paths:

\begin{lstlisting}[language=C++]
if (!pnp_.ready()) return;
if (!imu_valid_) {
  RCLCPP_WARN_THROTTLE(get_logger(), *get_clock(), 2000,
    "Waiting for /micro_imu");
  return;
}
\end{lstlisting}

The \code{return} exits the current function early. This pattern is called a
guard clause. It keeps the rest of the function from running when required
inputs are missing.

\section{Loops}

A loop repeats work. This package often loops over detections or four armor
faces:

\begin{lstlisting}[language=C++]
for (const auto & det : detections) {
  // use det
}
\end{lstlisting}

\code{const auto & det} means: for each detection, call it \code{det}, do not
copy it, and do not modify it.

\section{ROS 2 Nodes}

A ROS 2 node is a running component that communicates with other nodes. The
\code{AutoAimNode} class inherits from \code{rclcpp::Node}:

\begin{lstlisting}[language=C++]
class AutoAimNode : public rclcpp::Node
\end{lstlisting}

This gives it ROS features such as parameters, subscriptions, publishers, and
logging.

\section{Topics And Messages}

A topic is a named channel. A message type defines what data is sent on that
channel.

\begin{longtable}{p{0.32\textwidth}p{0.32\textwidth}p{0.26\textwidth}}
\toprule
Topic & Message type & Direction \\
\midrule
\topic{/detector/armors} & \code{vision_msgs/Detection2DArray} & input \\
\topic{/camera_info} & \code{sensor_msgs/CameraInfo} & input \\
\topic{/micro_imu} & \code{std_msgs/Float32MultiArray} & input \\
\topic{/tracker/cmd_gimbal} & \code{geometry_msgs/Twist} & output \\
\topic{/cmd_vel_AI} & \code{geometry_msgs/Twist} & output legacy \\
\topic{/tracker/aim_pixels} & \code{geometry_msgs/Twist} & output debug \\
\topic{/tracker/marker} & \code{visualization_msgs/MarkerArray} & output debug \\
\topic{/auto_aim/debug} & \code{auto_aim/AutoAimDebug} & output debug \\
\bottomrule
\end{longtable}

\section{Subscriptions}

A subscription tells ROS: when a message arrives on this topic, call this
function.

\begin{lstlisting}[language=C++]
det_sub_ = create_subscription<vision_msgs::msg::Detection2DArray>(
  "/detector/armors", rclcpp::SensorDataQoS(),
  std::bind(&AutoAimNode::detectionCallback, this, std::placeholders::_1));
\end{lstlisting}

This means \code{detectionCallback} runs once for each incoming detection
frame.

\section{Publishers}

A publisher sends messages:

\begin{lstlisting}[language=C++]
cmd_pub_ = create_publisher<geometry_msgs::msg::Twist>(
  "/tracker/cmd_gimbal", rclcpp::SensorDataQoS());

geometry_msgs::msg::Twist cmd;
cmd.angular.z = yaw_target_micro;
cmd.angular.y = pitch_target_micro;
cmd.angular.x = fire_decision ? 1.0 : 0.0;
cmd_pub_->publish(cmd);
\end{lstlisting}

The firmware contract depends on these exact fields. Do not move yaw, pitch,
fire, or distance to different fields without a matching firmware change.

\section{Parameters}

Parameters are runtime settings. The node declares them at startup:

\begin{lstlisting}[language=C++]
cfg.bullet_speed = declare_parameter("bullet_speed", 25.0);
\end{lstlisting}

If a YAML file provides \param{bullet_speed}, ROS uses the YAML value. If not,
the default \code{25.0} is used.

Robot-specific values belong in \file{src/config/params_realsense_16.yaml} or
\file{src/config/params_zed_64.yaml}, not hard-coded into C++.

\section{Launch Files}

Launch files start nodes and provide parameters. For example,
\file{auto_aim_realsense_16.launch.py} loads
\file{config/params_realsense_16.yaml}. Use launch files for normal robot runs
because they keep the startup command consistent.

\section{Logging}

ROS logging prints useful runtime information:

\begin{lstlisting}[language=C++]
RCLCPP_INFO(get_logger(), "Auto-aim node started");
RCLCPP_WARN_THROTTLE(get_logger(), *get_clock(), 1000, "PnP FAILED");
\end{lstlisting}

Throttled logs print at most once per time window. This prevents the terminal
from being flooded at 30 frames per second.

\section{External Libraries Used Here}

\begin{longtable}{p{0.22\textwidth}p{0.68\textwidth}}
\toprule
Library & Why this package uses it \\
\midrule
ROS 2 \code{rclcpp} & Nodes, parameters, publishers, subscribers, logging. \\
OpenCV & PnP, camera matrices, Rodrigues rotations, point projection. \\
Eigen & Linear algebra for EKF state, matrices, covariance, and Kalman math. \\
\code{tf2} & Quaternions, vector rotations, marker orientations. \\
\code{angles} & Correct wrap-around math for radians near $+\pi$ and $-\pi$. \\
\bottomrule
\end{longtable}

\section{How To Read A New File}

Use this routine when opening a file for the first time:

\begin{enumerate}
  \item Read the matching header first if one exists.
  \item Identify the main data types: structs and classes.
  \item Find the public functions. These are the file's interface.
  \item Find where those functions are called with \code{rg "functionName"}.
  \item Read implementation from the top-level function inward.
  \item Track units: meters, radians, pixels, seconds.
  \item Track frame: camera, gimbal body, or odom.
\end{enumerate}

\section{How To Make A Small Change Safely}

\begin{enumerate}
  \item State the behavior you want to change in one sentence.
  \item Find the earliest pipeline stage responsible for that behavior.
  \item Change only that module when possible.
  \item Rebuild with \code{colcon build --packages-select auto_aim}.
  \item Relaunch and record debug data.
  \item Compare \topic{/auto_aim/debug} before and after.
\end{enumerate}

Avoid changing several unrelated parameters at once. If the behavior changes,
you will not know which change caused it.

\section{Useful Terminal Commands}

\begin{longtable}{p{0.40\textwidth}p{0.50\textwidth}}
\toprule
Command & Purpose \\
\midrule
\code{rg "text" src} & Search the source tree quickly. \\
\code{rg --files src} & List source files. \\
\code{colcon build --packages-select auto_aim} & Build only this package. \\
\code{source install/setup.bash} & Load the built package into the current shell. \\
\code{ros2 topic list} & Show active ROS topics. \\
\code{ros2 topic echo /auto_aim/debug} & Print debug messages. \\
\code{ros2 topic hz /detector/armors} & Measure detector publish rate. \\
\code{ros2 param list /auto_aim} & Show parameters for the node. \\
\code{ros2 bag record ...} & Record data for offline debugging. \\
\bottomrule
\end{longtable}

\section{Common Beginner Mistakes}

\begin{itemize}
  \item Editing generated files under \file{build/} or \file{install/}. Edit
  \file{src/} instead.
  \item Forgetting to rebuild after editing C++ or message files.
  \item Forgetting to run \code{source install/setup.bash} after a build.
  \item Mixing degrees and radians.
  \item Mixing camera-frame and odom-frame coordinates.
  \item Hiding calibration mistakes with \param{pitch_offset_deg} or
  \param{yaw_offset_deg}.
  \item Changing the \topic{/tracker/cmd_gimbal} field mapping without a
  firmware update.
\end{itemize}

\section{First Exercises}

\begin{enumerate}
  \item Find where \param{bullet_speed} is declared in the C++ node.
  \item Find where \topic{/tracker/cmd_gimbal} is created and where messages
  are published on it.
  \item Find all files that mention \code{FireGate}.
  \item Change only a YAML value, relaunch, and confirm the parameter changed
  with \code{ros2 param get}.
  \item Record \topic{/auto_aim/debug} while pointing at a static target and
  identify which fields belong to detection, PnP, EKF, aim, command, and fire
  gate.
\end{enumerate}
